%% =====================================================================
%% Self-Propelled Arm System (S.P.A.S.) - SHORT REPORT, VERSION 2
%% Project report - Embedded Systems, Cyber-Physical Systems and Robotics
%% Technical University of Munich, Campus Heilbronn
%%
%% Scope: the live production runtime only. The monocular depth network
%% and everything built on it (depth stop criteria, depth obstacle
%% avoidance, depth grasp verification, the MiDaS development scripts and
%% the ArUco triangulation study) are deliberately out of scope, because
%% none of them is on the control path in the delivered configuration and
%% the demonstrated end-to-end mission ran with DepthNet disabled.
%%
%% Compile with:  pdflatex SPAS_technical_report_short_version2.tex (twice)
%%
%% Figures use \projectfig{path}{width}{height}{text}: the real file if it
%% exists, otherwise a labelled placeholder box, so this always compiles.
%% All figures in this version resolve to real files: the platform,
%% target, detection, grasp and release photographs were captured from
%% the recorded hardware demonstration, the state graph and the training
%% curves are project assets.
%% =====================================================================

\documentclass[conference,a4paper,10pt]{IEEEtran}

\usepackage[T1]{fontenc}
\usepackage[utf8]{inputenc}
\usepackage{graphicx}
\usepackage{xcolor}
\usepackage{booktabs}
\usepackage{amsmath}
\usepackage{url}
\usepackage{tabularx}
\usepackage[strings]{underscore}

\renewcommand{\arraystretch}{1.1}

\newcommand{\phbox}[4]{%
  \setlength{\fboxrule}{0.9pt}%
  \fcolorbox{black!45}{black!6}{%
    \begin{minipage}[c][#2][c]{\dimexpr#1-2\fboxsep-2\fboxrule\relax}
      \centering\scriptsize
      \textbf{[\,PLACEHOLDER IMAGE\,]}\\[3pt]
      \itshape #3\\[4pt]
      \upshape\tiny\ttfamily #4
    \end{minipage}}%
}

\newcommand{\projectfig}[4]{%
  \IfFileExists{#1}{\includegraphics[width=#2]{#1}}{\phbox{#2}{#3}{#4}{#1}}%
}

\begin{document}

\title{Self-Propelled Arm System (S.P.A.S.):\\
An Autonomous Can-Collecting Robot\\
on a Monocular Tracked Platform}

\author{%
  \IEEEauthorblockN{Filip Chekalkin \quad Huaijing Hou \quad Spaska Janeva \quad Kamen Kanev\\
                    Giorgi Khundadze \quad Xiaosheng Tan \quad Stoyan Vlaev \quad Weiye Yuan}
  \IEEEauthorblockA{Embedded Systems, Cyber-Physical Systems and Robotics\\
                    Technical University of Munich, Campus Heilbronn, Germany}
}

\maketitle

%% The abstract was removed by request. It is kept here, commented out, in
%% case it is ever wanted back: uncomment the block below.
%
% \begin{abstract}
% This report describes the Self-Propelled Arm System (S.P.A.S.), an
% autonomous robot that searches an indoor area for aluminium drink cans,
% drives to them, grasps them with a four-degree-of-freedom arm, and
% deposits them into a marked bin. The platform is a tracked Hiwonder
% JetTank with an NVIDIA Jetson Nano and a single $320\times240$ RGB
% camera. It has no depth sensor, no range finder and no wheel encoders, so
% the base is a pure actuator and every control decision is closed on image
% appearance. The report is scoped to the production runtime: the deployed
% single-stage can detector, the AprilTag bin detector, the mission state
% machine, the visual servoing stages and the grasp sequence. The deployed
% detector reaches an $F_1$ score of $0.924$ on robot video, and a complete
% search-grasp-transport-release mission was executed on hardware in
% $73.8$~s with the depth network switched off. We report the measured
% platform constraints that bound the design and one silent failure mode
% found by inspecting the code rather than by testing.
% \end{abstract}

%% ---------------------------------------------------------------------
\section{Problem Statement}
%% ---------------------------------------------------------------------

Indoor spaces accumulate small litter, and drink cans are among the most
common items. Collecting them is repetitive, low-risk work and therefore
a reasonable target for a small mobile manipulator. It is also a useful
teaching problem because the robot must solve perception, navigation,
decision-making, and manipulation simultaneously, and any failure in one
is immediately visible.

The mission is: search for a can, drive to it while keeping it in view,
stop within reach, grasp and lift it, find a bin carrying a fiducial
marker, drive to the bin, release the can, and either resume searching or
report completion.

What makes the problem interesting is the hardware budget. The platform
carries one low-resolution RGB camera and nothing else: no stereo pair,
no depth camera, no range finder, no inertial unit the software reads,
and no wheel encoders. The tracked base accepts a direction and a
normalised speed and returns no measurement at all. The on-board computer
is a Jetson Nano with 4~GB of memory shared between processor and
graphics unit.

\textbf{Scope of this report.} The project also built a monocular depth
capability and several exploratory tools around it. None of that is on
the control path of the delivered system. In the shipped configuration
both approach stop criteria are appearance-based, obstacle avoidance is
disabled, and grasp verification is disabled, so nothing decisive
consults depth; the depth network's only remaining consumer is coarse map
bookkeeping that by construction cannot trigger an action. The archived
end-to-end mission analysed in Section~\ref{sec:results} in fact ran with
the depth network disabled outright. This report therefore describes only
the live runtime, and every claim below concerns a purely
appearance-based control loop. The complete implementation is available
in the project repository~\cite{spas_repo}.

%% ---------------------------------------------------------------------
\section{System Overview}
%% ---------------------------------------------------------------------

The robot is built on a Hiwonder JetTank~\cite{hiwonder}, a tracked
chassis carrying a Jetson Nano, a serial-bus servo arm and a
forward-facing camera. Table~\ref{tab:hardware} lists the components, and
Fig.~\ref{fig:hw} shows the platform and the two objects it has to
recognise.

\begin{table}[t]
\caption{Hardware components}
\label{tab:hardware}
\centering
\footnotesize
\begin{tabularx}{\columnwidth}{@{}lX@{}}
\toprule
\textbf{Component} & \textbf{Function} \\
\midrule
JetTank chassis & Tracked base. Accepts normalised speed commands only;
returns no measurement. \\
Jetson Nano (4~GB) & Runs all perception, planning and control. \\
RGB camera & The only exteroceptive sensor, $320\times240$ pixels,
horizontal field of view $\approx 60^{\circ}$. \\
Arm and gripper & Five serial-bus servos over a TTL chain. Position is
readable; grip force is not. \\
AprilTag & 80\,mm \texttt{tag36h11}, identifier 0, mounted on the bin. \\
\bottomrule
\end{tabularx}
\end{table}

The most useful way to characterise this hardware is by what it reports
back, because that asymmetry determines the software design.

\textbf{The base is a pure actuator.} No encoders, no current sensing, no
inertial measurement. Every base parameter is an empirical constant
obtained by driving the robot and recording the results, and nothing in
the running system can verify those constants.

\textbf{The camera is the only sensor.} Every observation about the
outside world passes through one device. If it is saturated, blurred or
pointed the wrong way, the robot is blind, and the design must survive
that rather than prevent it.

\textbf{The gripper is force-blind.} The servos accept an angle and a
speed and can report position, which is the one place the system has
genuine proprioception. They do not expose grip force or motor current,
so the gripper closes to a commanded angle without knowing whether
anything is inside it.

The software exploits the one asymmetry available: arm motions are
verified by polling servo positions until they stop changing, whereas
base motions can only be verified by looking at the camera again.

\begin{figure}[t]
\centering
\begin{minipage}[b]{0.485\columnwidth}
\centering
\projectfig{figures/robot_platform.jpg}{\linewidth}{3.1cm}%
  {Photograph of the assembled robot: tracked base, arm with gripper,
   forward-facing camera.}
\end{minipage}\hfill
\begin{minipage}[b]{0.485\columnwidth}
\centering
\projectfig{figures/targets_and_bin.jpg}{\linewidth}{3.1cm}%
  {Photograph of the 250\,ml can and the bin carrying the printed
   AprilTag.}
\end{minipage}
\caption{The platform carrying a grasped can (left) and the two target
objects in the arena (right). The can is a 250\,ml slim aluminium can,
53\,mm in diameter and 133\,mm tall, weighing about 11\,g empty. The bin
is a cardboard box roughly 5\,cm tall with the printed
\texttt{tag36h11} marker on its side. Both photographs are frames from
the recorded hardware demonstration.}
\label{fig:hw}
\end{figure}

%% ---------------------------------------------------------------------
\section{Perception on the Control Path}
%% ---------------------------------------------------------------------

Two detectors run on the robot, and between them they supply every
measurement the controller uses.

\textbf{Can detection} uses an SSD-MobileNet-V1 model~\cite{ssd,
mobilenet} at $300\times300$ input resolution with a single foreground
class, exported to ONNX and executed through the
\texttt{jetson-inference} \texttt{detectNet} interface~\cite{jetsoninference}
at a confidence threshold of $0.20$ and a clustering overlap of $0.30$.
This architecture was chosen over the YOLO model~\cite{yolo,ultralytics}
used for offline dataset work because it is the path NVIDIA supports
natively on the Nano; the trade is quantified in
Section~\ref{sec:results}. The detector returns a bounding box, from
which the controller derives two scalars: the normalised horizontal error
and the normalised box height. Fig.~\ref{fig:canpredict} shows both the
raw detector output and the same two scalars as they appear in the
on-board telemetry overlay during an approach.

\textbf{Bin detection} locates an 80~mm \texttt{tag36h11}
AprilTag~\cite{apriltag} with identifier 0 using the
\texttt{pupil\_apriltags} library, with OpenCV's ArUco module as a
fallback backend. Beyond the identifier, the detector computes three
descriptors from the four corners, and these are what the docking
controller consumes:

\begin{itemize}\itemsep2pt
  \item the \emph{normalised tag height}, which grows as the robot
        approaches and serves as the distance proxy;
  \item the \emph{vertical edge error}, the relative difference between
        the projected lengths of the left and right edges,
        $e_{\mathrm{v}} = (\ell_{r} - \ell_{l}) /
        \tfrac{1}{2}(\ell_{r} + \ell_{l})$, which is zero when the camera
        faces the tag squarely and changes sign with the viewing angle;
  \item the \emph{maximum corner angle error}, the largest deviation of
        an interior corner from $90^{\circ}$, which gates how much the
        other two should be trusted.
\end{itemize}

A metric pose is available by solving the perspective-n-point problem if
a calibration file is supplied, but the delivered controller does not use
it. The three descriptors above need no calibration, because the focal
length cancels in the ratio, and calibration is one more artefact that
can go stale. Fig.~\ref{fig:tag} shows the tag detection at close range,
where the docking controller reads a tag height of $0.186$ and a
horizontal error of $-0.055$ from the same four corners.

\begin{figure}[t]
\centering
\begin{minipage}[b]{0.36\columnwidth}
\centering
\projectfig{../can_predict.png}{\linewidth}{2.6cm}%
  {A camera frame with the detected can and its confidence score.}
\end{minipage}\hfill
\begin{minipage}[b]{0.60\columnwidth}
\centering
\projectfig{figures/onboard_can_detection.jpg}{\linewidth}{2.6cm}%
  {On-board frame during the approach stage with the telemetry overlay.}
\end{minipage}
\caption{Can detection. Left: the detector's own output on a
$300\times300$ frame, with the confidence score. Right: the same
detection inside the running mission, with the telemetry overlay
reporting the state (\texttt{APPROACHING}), the commanded base motion,
and the two scalars the controller actually uses,
$e_x = -0.119$ and a normalised box height of $0.255$. The entire
perception stack operates at $320\times240$.}
\label{fig:canpredict}
\end{figure}

\begin{figure}[t]
\centering
\projectfig{figures/onboard_tag_detection.jpg}{0.86\columnwidth}{4.4cm}%
  {On-board frame during bin docking with the AprilTag detected on the
   side of the bin.}
\caption{AprilTag detection during bin docking. The green box marks the
detected \texttt{tag36h11} marker; the overlay reports a tag height of
$0.186$ against the $0.15$ stop threshold and a horizontal error of
$-0.055$. Note also \texttt{grabbed=True} and \texttt{DEPTH=NO}: the run
carried a can and used no depth information.}
\label{fig:tag}
\end{figure}

%% ---------------------------------------------------------------------
\section{Mission Control}
%% ---------------------------------------------------------------------

The mission is executed by an explicit event-driven finite state machine
with fifteen states, shown in Fig.~\ref{fig:fsm}. Initialisation starts
the camera and parks the arm. Planning selects the next target: a
remembered can, a fresh search, a patrol leg, the bin, or completion.
Searching, aligning, approaching and final verification are the four
visual stages described in Section~\ref{sec:nav}. Finalising runs either
the grasp sequence or the release sequence. An intermediate state absorbs
timeouts, and two terminal states end the run.

Transitions are declared as an explicit table of (state, event) pairs,
and any pair not in the table raises an error rather than silently doing
nothing. This is deliberate: an undeclared transition is a specification
bug, and it is far easier to find when the program stops than when the
robot quietly continues in an unexpected state. It also makes the
transition relation a small finite object that can be reasoned about
directly. Two useful properties follow from reading it. The finalising
state has exactly two incoming edges, both requiring the target-stable
event, so a grasp can only follow several consecutive well-aligned
frames. The map-navigation state has four outgoing edges, and none of
them reaches finalising, so the dead-reckoned position can never, by
itself, trigger a physical action.

Failure handling is uniform. Any state that exceeds its time or step
budget emits a timeout, which is routed to the intermediate state; that
state increments a retry count and either replans or, once the limit is
exhausted, fails. If something goes wrong unexpectedly, the system
switches to a failure state. No matter what happens, the run function
always handles cleanup in a \texttt{finally} block: the base is stopped,
the arm moves back to its safe position, and the camera is released.
Cleanup can safely run more than once, thanks to a lock that prevents any
issues if it gets called repeatedly during an interrupted mission.

All tunable values are stored in two JSON files, organised by decision
type instead of by module. One file covers topics such as which paths to
take, simulation flags for the hardware, and how many times the system
should retry if something fails. The other file is where all the values
discovered from trial and error go---speeds, timeouts, tolerances,
thresholds, and the different arm positions. When it is time to use these
settings, both files are combined in a set order, and any command-line
changes take priority at the end. Speeds are referenced symbolically, for
example, \texttt{turn.slow}, and resolved against a small profile table,
so re-tuning the platform does not mean editing dozens of call sites.

Each hardware device has an independent simulation flag, and all three
default to simulated, so the complete state machine can be exercised on
a laptop without a robot attached. On the robot, devices are enabled one
at a time, with the arm last. This is the single most effective safety
measure in the project because the arm is the component most likely to
damage itself.

\begin{figure*}[t]
\centering
\projectfig{../writing_assets/visualization/abstract_v1.png}{0.78\textwidth}{7.5cm}%
  {Design-level mission state graph with guard conditions on the
   transitions and the shared mission variables listed in the top-left
   box.}
\caption{The mission state machine. Guards such as \texttt{target\_found}
and \texttt{grabbed} are fields of a shared mission context. The
implementation refines this graph, notably by splitting approach into a
coarse stage and a final verification stage, but preserves its
structure.}
\label{fig:fsm}
\end{figure*}

%% ---------------------------------------------------------------------
\section{Navigation and Manipulation}
\label{sec:nav}
%% ---------------------------------------------------------------------

\subsection{The Four Visual Stages}

Can navigation and bin navigation share four stages with different
parameters and detectors. For a detection with centre $x_c$ in an image
of width $W$, the normalised horizontal error is

\begin{equation}
  e_x = \frac{x_c - W/2}{W} \in [-0.5,\ 0.5],
  \label{eq:errorx}
\end{equation}

negative when the target lies left of centre. This single scalar drives
every steering decision in the system.

During the \textbf{searching} phase, the system rotates in a designated
direction until either the detector identifies a target or the state
times out. In the \textbf{aligning} phase, the base turns toward the
target until the absolute horizontal error ($|e_x|$) is within a
specified tolerance. If the target is lost during this process, the base
halts immediately, and the lost-frame counter is incremented. When this
counter reaches a predefined threshold, the stage concludes by reporting
the target as missing, preventing unnecessary rotation without a target.

Stopping immediately on loss matters, because a rotating robot that has
lost its target will quickly turn past it. \textbf{Approaching}
re-detects every frame, steers when $|e_x|$ exceeds a tolerance and
otherwise drives forward, stopping when the target's normalised
bounding-box height exceeds a threshold: $0.36$ for the can and $0.15$
for the bin. Apparent size is used as the distance proxy because it
depends only on the detector. \textbf{Final verification} repeats
alignment with a tighter tolerance and requires several consecutive
frames to pass before the mission commits.

Because the base returns nothing, the system also keeps a deliberately
approximate spatial memory, integrating the commands it issued to
estimate its own pose and to remember roughly where cans were seen. Its
purpose is attention, not localisation: a visible target always overrides
it, and arriving at a remembered coordinate only triggers a visual search
there. Section~\ref{sec:results} explains why nothing stronger would be
defensible.

\subsection{Grasping}

The arm follows a sequence of predefined positions, in place of real-time
inverse kinematics computation. Since the task parameters remain
consistent---the same object, a flat surface, and a fixed approach
distance---a general kinematic solution would increase complexity without
offering substantial advantages for this specific scenario. The sequence
is: pause for the base to settle; move to \texttt{arm\_down} and wait for
the servos to actually arrive; drive the base slowly forward for a fixed
time, pushing the can into the open gripper; hold still; close to
\texttt{grab}; lift to \texttt{carry}. Fig.~\ref{fig:grasp} shows the
four decisive moments of that sequence as the robot saw them.

Two steps deserve comment. The settle wait is the one place with real
feedback: rather than sleeping for a fixed duration, the controller polls
each servo's reported position until consecutive readings differ by less
than a small threshold for several samples, with a hard timeout. If the
arm has not settled, the base push is blocked outright.

If the robot moves forward before the gripper has finished lowering, it
risks pushing the can away instead of picking it up. This technique,
pushing the can into the gripper rather than closing the gripper around a
stationary can, is somewhat unconventional, but it suits the hardware's
constraints. The arm's reach is short, and the gripper's closure is
somewhat imprecise, so the robot first lowers the open gripper and then
moves forward to let the can slide inside. This approach avoids the need
for fine manipulation and instead makes use of the base's more reliable
movement.

The cost is that it works only for a light object on a low-friction
floor.

\begin{figure}[t]
\centering
\projectfig{figures/grasp_sequence.jpg}{\columnwidth}{6.0cm}%
  {Four on-board frames covering final verification, the base push, the
   gripper closing and the lift.}
\caption{The grasp sequence seen from the robot's own camera, reading
left to right and top to bottom. Top left: final verification, with the
can aligned to $e_x = -0.063$ and a box height of $0.407$ past the
$0.36$ threshold. Top right: the base pushing forward at speed $0.20$
with the gripper already lowered. Bottom left: the gripper closing on the
can. Bottom right: the can lifted to the carry pose. The overlay's
\texttt{DEPTH=NO} field confirms no depth information was used.}
\label{fig:grasp}
\end{figure}

%% ---------------------------------------------------------------------
\section{Results}
\label{sec:results}
%% ---------------------------------------------------------------------

\subsection{Detector Accuracy}

The deployed SSD-MobileNet-V1 model was evaluated against the custom
model it replaced, on frozen photo and video sets captured by the robot.
It reaches $F_1 = 0.839$ on photographs and $0.924$ on video at its
$0.20$ operating threshold, against $0.899$ and $0.945$ for the previous
model at its own $0.50$ threshold. The deployed model is therefore
measurably weaker, by about six points on photographs and two on video,
and was still chosen because it removes a custom post-processing layer
and a custom inference backend from the robot in favour of the
vendor-supported path. On a fixed deadline, a slightly less accurate
detector that is reliable to build and deploy is worth more than a better
one that is fragile. The ONNX export was verified numerically: outputs
differed from the original network by at most $4.4\times10^{-7}$, so the
gap is a property of the architecture, not the conversion.

For reference, the offline YOLO11n model used for dataset iteration
achieved a mean average precision of $0.995$ at an
intersection-over-union threshold of $0.5$ on its validation split, with
the training behaviour shown in Fig.~\ref{fig:training}. That figure
should be read sceptically: the dataset is small, single-class and
drawn from one room, so it describes how easy this detection problem is
in this arena rather than how good the detector is.

\begin{figure}[t]
\centering
\projectfig{figures/training_curves.pdf}{\columnwidth}{4.2cm}%
  {Training and validation loss curves and validation mAP for the
   offline YOLO11n model.}
\caption{Offline YOLO11n training used for dataset iteration only; this
model does not run on the robot. Validation mAP@0.5 saturates near
$0.995$ within about twenty epochs, which reflects the narrowness of the
single-room, single-class dataset more than the difficulty of the
detection problem in general.}
\label{fig:training}
\end{figure}

\subsection{End-to-End Mission}

Table~\ref{tab:trace} reconstructs an archived successful mission
executed on hardware, taken from its log. The run used only the camera,
with the depth network disabled. Fig.~\ref{fig:release} shows the final
moment of such a run on the physical platform.

\begin{table}[t]
\caption{Reconstruction of a complete successful mission}
\label{tab:trace}
\centering
\footnotesize
\begin{tabularx}{\columnwidth}{@{}lX@{}}
\toprule
\textbf{Phase} & \textbf{Observed behaviour} \\
\midrule
Search can   & Accepted on the first evaluated frame at confidence
$0.586$, $e_x = -0.388$. \\
Align can    & 72 corrective turn steps to bring $|e_x|$ inside the
$0.15$ tolerance. \\
Approach can & 10 forward pulses; box height grew $0.176 \to 0.509$
against a $0.400$ stop threshold. \\
Near align   & 15 further steps at the tighter tolerance. \\
Grasp        & Arm settled in $4.94$\,s; $7.0$\,s push at speed $0.1$;
$5.0$\,s settle; grip and lift. \\
Bin phases   & Tag acquired, 9 alignment steps, 35 forward pulses to a
tag height of $0.590$. \\
Release      & Gripper opened, arm returned to its safe pose. \\
\midrule
\textbf{Total} & \textbf{73.8\,s, mission reported successful} \\
\bottomrule
\end{tabularx}
\end{table}

\begin{figure}[t]
\centering
\projectfig{figures/bin_release.jpg}{\columnwidth}{4.4cm}%
  {Photograph of the robot holding the grasped can over the bin
   immediately before release.}
\caption{The end of a complete mission on hardware: the robot has docked
against the bin and holds the grasped can over it, immediately before the
gripper opens and the arm returns to its safe pose.}
\label{fig:release}
\end{figure}

Across the run, the can detector was invoked 98 times and the tag
detector 41 times, with confidence scores from $0.558$ to $0.942$. Three
observations follow. Alignment dominates the time budget: 72 steps were
needed to remove an initial error of $-0.388$, because each corrective
pulse is deliberately small. The visual stop criterion behaved as
designed, with box height growing almost monotonically and the stage
halting on the first frame past threshold; the one non-monotonic step, a
dip from $0.204$ to $0.196$, is the detector jitter such a threshold must
tolerate. Roughly twelve seconds of the run are spent deliberately not
moving, in the settle wait, the slow push and the post-push lock, which
is the price of grasping reliably with imprecise hardware.

An archived failed run is equally instructive: it terminated after
$19.6$\,s when the can was lost for more consecutive frames than the
limit allows. The robot stopped cleanly rather than driving blindly, as
the design intends.

\subsection{Measured Constraints}

Table~\ref{tab:constraints} lists the platform limits measured during
development. These were the binding constraints on the design.

\begin{table}[t]
\caption{Measured platform constraints}
\label{tab:constraints}
\centering
\footnotesize
\begin{tabularx}{\columnwidth}{@{}Xl@{}}
\toprule
\textbf{Constraint} & \textbf{Value} \\
\midrule
Turn rate asymmetry, left versus right, at identical commands
& 1--3\,rad/s \\
Usable range of the on-board can detector & $\approx 3$\,m \\
Reliable AprilTag localisation range & $\approx 1$\,m \\
Sustained processing rate & 10--15\,frames/s \\
Maximum gripper reach height & $\approx 20$\,cm \\
\bottomrule
\end{tabularx}
\end{table}

The turn asymmetry is the most damaging entry, and its magnitude is worth
spelling out. The dead-reckoning model uses a single constant of $\pi$
radians per speed-second, so the configured turn profiles predict
$0.471$~rad/s slow and $0.942$~rad/s fast. A direction-dependent
discrepancy of $1$ to $3$~rad/s is therefore between $1.1$ and $6.4$
times the modelled rate itself. This is stronger than saying the constant
needs calibrating: a single scalar cannot describe a platform whose
behaviour differs between two directions by more than the modelled
magnitude, so the model has the wrong form. Nothing that depends on
integrated heading can be trusted, which is precisely why every decisive
loop in the system is closed on image appearance and why the spatial
memory is never allowed to trigger an action.

\subsection{A Silent Failure Mode}

Reading the code against the mission's intent surfaced one defect that
testing had not. The grasp verification step is guarded by a
configuration flag that is set to false in the shipped configuration,
so the pickup is recorded unconditionally. If the push fails to capture
the can, the mission context still asserts that an object is held, the
robot drives to the bin, performs the release over it, and reports
success with an empty gripper. Nothing downstream can detect this,
because no later state re-examines the gripper.

This is the only one of the system's failure modes that is silent. Target
loss, state timeouts, a stalled arm, a stale map entry and an
unavailable camera frame all produce an explicit failure or a recovery.
A system that fails loudly in five ways and silently in one should spend
its next increment of effort on the silent one. The property is restored
by enabling the flag; the deeper fix is gripper force or current sensing,
which this hardware does not expose.

%% ---------------------------------------------------------------------
\section{Limitations and Conclusion}
%% ---------------------------------------------------------------------

The system does what it set out to do, but the honest description is
narrow. The pose estimate integrates commanded motion without correction,
so its error is unbounded and its model is invalid given the measurement
above. The detector was trained and validated on data from a single room
with a single object type, so its metrics do not generalise. The
push-based grasp depends on floor friction that was never measured, and
by the finding above, the system does not notice when that assumption
fails. At 10--15 frames per second, the visual control loop is slow
relative to the base, which is why motion is pulsed, and speeds are low.

Taking these limitations into account, the core design choice can be
understood as a practical one. Since each decisive loop relies on image
measurements rather than estimated geometry, most errors tend to degrade
performance gradually rather than directly lead to failure. For example,
if the map drifts, the robot needs to search a bit more; if a turn is
asymmetric, it just requires some extra correction steps. With the single
exception of the empty-gripper case, the system is slow and fussy but
rarely confidently wrong, and on hardware with no proprioception, that
is the property worth engineering for.

The contributions are the integration and the choices it forced. An
explicit state machine with a declared transition table and uniform
timeout handling made a fifteen-state mission debuggable and made two
useful safety properties readable straight off the table.
Appearance-based visual servoing allowed the system to tolerate a base
with no odometry and a turn model that does not hold. The system uses an
intentionally approximate spatial memory, allowing the robot to keep
track of can locations without assuming it knows its precise position.
With a push-based grasp, a challenging precision task is replaced with a
simpler one that suits the hardware's capabilities.

The next steps for improvement are relatively inexpensive. For example,
enabling grasp verification can address the silent failure mode at the
cost of just one additional observation per pickup. Calibrating separate
left- and right-turn constants, which the configuration already supports
and currently leaves at unity, would address the largest single source of
motion error. Adding one inexpensive forward-facing distance sensor would
give the system its first true metric measurement, and gripper force
sensing would be the proper fix for the failure mode above.

%% ---------------------------------------------------------------------
\begin{thebibliography}{9}
\footnotesize

\bibitem{spas_repo}
K.~Kanev \emph{et al.}, ``Self-Propelled Arm System (S.P.A.S.),'' project
repository, 2026. [Online]. Available:
\url{https://github.com/kamkanev/Self-Propelled-Arm-System}

\bibitem{yolo}
J.~Redmon, S.~Divvala, R.~Girshick, and A.~Farhadi, ``You Only Look Once:
Unified, Real-Time Object Detection,'' in \emph{Proc. CVPR}, 2016,
pp.~779--788.

\bibitem{ultralytics}
G.~Jocher, A.~Chaurasia, and J.~Qiu, ``Ultralytics YOLO,'' 2024.

\bibitem{ssd}
W.~Liu \emph{et al.}, ``SSD: Single Shot MultiBox Detector,'' in
\emph{Proc. ECCV}, 2016, pp.~21--37.

\bibitem{mobilenet}
A.~G.~Howard \emph{et al.}, ``MobileNets: Efficient Convolutional Neural
Networks for Mobile Vision Applications,'' arXiv:1704.04861, 2017.

\bibitem{apriltag}
E.~Olson, ``AprilTag: A Robust and Flexible Visual Fiducial System,'' in
\emph{Proc. ICRA}, 2011, pp.~3400--3407.

\bibitem{jetsoninference}
D.~Franklin, ``jetson-inference: Deploying Deep Learning with NVIDIA
TensorRT,'' NVIDIA, 2024.

\bibitem{hiwonder}
Hiwonder, ``JetTank ROS Robot Tank with Jetson Nano,'' product
documentation, 2024.

\end{thebibliography}

\end{document}
